\section*{Заключение}
\addcontentsline{toc}{section}{Заключение}

В результате выполнения курсового проекта было разработано веб-приложение для хранения, систематизации и публикации учебных материалов с ролевым разграничением доступа. В ходе работы была проанализирована предметная область, рассмотрены существующие подходы к организации электронных архивов учебных материалов, а также выбран технологический стек, включающий Python, Flask, Flask-Login, Flask-SQLAlchemy, SQLite, Jinja2 и Bootstrap.

На этапе проектирования были определены роли пользователей, основные сценарии использования системы, структура интерфейса и модель данных для сущностей, связанных с хранением материалов и файлов. На этапе разработки была реализована базовая структура приложения, аутентификация и авторизация пользователей, проверка прав доступа в соответствии с ролью, CRUD-интерфейс для пользователей, категорий и материалов, механизм загрузки, скачивания и удаления файлов, а также страница статистики с агрегированными данными.

Проведенное тестирование основных пользовательских сценариев подтвердило корректность работы реализованного приложения. Система обеспечивает доступ к функциям в соответствии с ролью пользователя, поддерживает работу с материалами и файлами и удовлетворяет обязательным требованиям курсового проекта. Полученный результат может рассматриваться как работоспособная основа для дальнейшего расширения системы, например для добавления более сложного поиска, журналирования действий пользователей, расширенной аналитики и интеграции с иными образовательными сервисами.